\section{Diskussion und Reflexion}\label{sec:diskussion}

\subsection{Beantwortung der Forschungsfragen}\label{sec:antworten}

\paragraph{FF1, Erkennungsansatz.} Variante~C ist für die App die am besten
belegte Wahl. In der letzten Datenmengenreihe erreicht der Detektor bei
200'000 Trainingsbildern auf dem realen Validierungsset eine \map{50} von
0.908. Er liefert Kartenklasse und Position mit einem Modell und läuft in
den eingereichten Apps in Echtzeit. Variante~A liefert keine Position.
Bei Variante~B sind die Lokalisierung und die Klassifikation ideal
entzerrter Ausschnitte gemessen, aber nicht die gesamte Erkennungskette.
Zudem fehlt für die letzte Reihe ein Latenzvergleich aller Ansätze auf
demselben Gerät. Wie gross der Vorteil von C bei Genauigkeit und
Rechenaufwand tatsächlich ist, lässt sich daher nicht beziffern.

\paragraph{FF2, Übertragung auf reale Aufnahmen.} Das Modell C wurde nur
mit synthetischen Bildern trainiert und erkennt Karten auf den eigenen
realen Aufnahmen. Damit funktioniert die Übertragung für das vorhandene
Validierungsset. Im Projekt wurden die Trainingsbilder und die reale
Prüfbasis schrittweise um flachere Perspektiven, andere
Lichtverhältnisse und schwierigere Negativbeispiele erweitert. Diese
Änderungen wurden nicht einzeln unter gleichen Bedingungen gemessen.
Wie gut das Modell auf neuen Aufnahmesessions oder anderen
Kartendesigns arbeitet, bleibt offen.

\subsection{Grenzen der Ergebnisse}\label{sec:grenzen}

Das Validierungsset enthält 813 Frames, davon 513 mit gelabelter Karte.
Viele Frames stammen aus denselben Videos und ähneln sich. Ausserdem
wird die beste Trainingsepoche auf diesem Set ausgewählt und dort auch
bewertet. Ein unabhängiges Testset und Wiederholungen mit weiteren
Trainings-Seeds fehlen. Die Kennzahlen können deshalb zu günstig
ausfallen. Kleine Unterschiede zwischen den Modellen sind nicht als
gesicherte Effekte zu verstehen.

Die Datenmengenreihe verwendet für jede Bildzahl denselben Generator
und Trainingsplan. Die Datensätze sind aber keine ineinander
enthaltenen Teilmengen. Mit der Bildzahl ändert sich auch, welche
Winkel und Kartenblätter auf die einzelnen Bilder fallen. Der
beobachtete Verlauf zeigt, was die grösseren Datensätze dieser Reihe
leisten. Er trennt den Einfluss der reinen Bildanzahl nicht vollständig
vom Einfluss der gezogenen Bilder.

Die Modelle werden mit unterschiedlichen Kennzahlen beurteilt.
\map{50} für die Detektoren und Top-1-Genauigkeit für die
Klassifikatoren dürfen nicht als dieselbe Trefferquote gelesen
werden. Vor allem sagt keine dieser Kennzahlen direkt, wie oft die
App einen ganzen Stapel richtig zählt.

Auch die Trainingsbilder haben Grenzen. Die Kartenscans stammen aus je
einem physischen Deck pro Blatt. Der Generator verändert Perspektive,
Licht und Hintergrund, aber nicht das Kartendesign. Hände, gekippte
Karten, harte Schatten und Spiegelungen auf glänzenden Karten fehlen.
Für solche Situationen und andere Decks braucht es zusätzliche reale
Prüfaufnahmen.

\subsection{Kurskorrekturen im Projekt}\label{sec:verworfen}

Zuerst sollten die Trainingsbilder fotografiert und von Hand gelabelt
werden. Für die benötigte Menge war das nicht praktikabel. Der Wechsel
zu gerenderten Bildern machte grössere Datensätze und automatische
Labels möglich. Reale Fotos dienen nun als Quellmaterial und zur
Prüfung, nicht als Trainingsbilder für Kartenhaufen.

Generator und Validierungsset wurden während des Projekts erweitert.
Besonders Aufnahmen aus flacheren Winkeln und Bilder ohne eindeutig
lesbare Karte zeigten, welche Situationen noch fehlten. Daraufhin
wurden die Kamerawinkel und die Negativbeispiele breiter angelegt.
Diese Schritte gehören zur Entwicklung des Systems. Da sie nicht
einzeln geprüft wurden, lässt sich ihnen keine bestimmte
Leistungssteigerung zuordnen. Entscheidend ist, dass die reale
Prüfbasis nicht nur Fälle zeigt, die der Generator bereits gut
abbildet.

Auch die Einstellungen der App mussten genauer betrachtet werden.
Konfidenzschwelle und Stabilitätsregel waren zunächst gesetzt, ohne
dass ihre Wirkung an mehreren Sessions gemessen worden war. Bei der
Prüfung zeigte sich zudem, dass die Apps die Konfidenz zeitweise
unterschiedlich auswerteten. Solche Werte müssen nachvollziehbar
festgehalten und mit aufgezeichneten Zählvorgängen geprüft werden.

\subsection{Eigene Werkzeuge}\label{sec:eigene-werkzeuge}

Das Dataset Tool verbindet die Sichtkontrolle der gerenderten Bilder,
das Labeln realer Aufnahmen und die Analyse von Zählsessions. Die
Dataset-CLI erzeugt und prüft die Daten. Beide Werkzeuge wurden für
die Fragen dieses Projekts gebaut. Dadurch liess sich eine Änderung
am Generator rasch an Bildern und Labels kontrollieren. Ein Label
kann zudem für alle drei Erkennungsansätze verwendet werden. Für
grössere Teams, die gemeinsam labeln und Ergebnisse gegenseitig
prüfen, wäre ein etablierter Dienst weiterhin sinnvoll.

\subsection{Rolle der KI und persönlicher Beitrag}\label{sec:ki-reflexion}

Der Code des Projekts wurde überwiegend mit Claude Code erstellt. Der
Verfasser legte die Aufgabe, die drei Ansätze, die Datenregeln und die
Architektur fest. Er fotografierte die Karten und die realen
Prüfaufnahmen, labelte die Daten, kontrollierte die erzeugten Bilder
und prüfte die Apps. Auch die Auswahl der Auswertungen und die
Interpretation der Ergebnisse lagen bei ihm. Die genaue Aufteilung
ist in Anhang~\ref{app:ki-deklaration} dokumentiert.

Die Arbeit mit KI machte die Umsetzung von Generator, Werkzeugen,
Trainingspipeline und zwei Apps in der verfügbaren Zeit möglich. Sie
nahm dem Verfasser die Prüfung nicht ab. Verdrehte Scans und Labels
fielen erst bei der Sichtkontrolle auf. Unterschiede bei der
Konfidenzauswertung der Apps mussten anhand der tatsächlichen
Ausgaben geprüft werden. Die Verantwortung für Annahmen, Messungen
und Schlussfolgerungen bleibt beim Verfasser.

\subsection{Lehren aus der Arbeit}\label{sec:lessons}

Eine reale Prüfbasis muss früh entstehen und mit neuen Fehlerfällen
wachsen. Die Regeln für lesbare, verdeckte und verwischte Karten
sollten vor dem Labeln feststehen. Ein reproduzierbarer Generator
hilft beim Vergleichen, ersetzt aber kein unabhängiges Testset.
Ebenso wichtig ist die Prüfung der vollständigen Zählung in der App.
Erst damit lässt sich beurteilen, ob gute Kennzahlen für einzelne
Frames auch zu einer verlässlichen Nutzung am Spieltisch führen.
